% Ficha de estudio 7 del Taller virtual ATA EPE: mecánica analítica II (lagrangiano reducido y trabajo virtual).
\documentclass[11pt,a4paper]{article}
% Preámbulo común de las fichas de estudio del Taller virtual ATA EPE.
\usepackage[spanish,es-noshorthands,es-tabla]{babel}
\usepackage[utf8]{inputenc}
\usepackage[T1]{fontenc}
\usepackage{lmodern}
\usepackage[margin=2.2cm]{geometry}
\usepackage[table]{xcolor}
\usepackage{amsmath,amssymb,bm}
\usepackage{tabularx,array,booktabs}
\usepackage{enumitem}
\usepackage{fontawesome5}
\usepackage[most]{tcolorbox}
\usepackage[hidelinks]{hyperref}
\usepackage{listings}

\definecolor{morado}{HTML}{2A1B7E}
\definecolor{amarillo}{HTML}{FFCF2E}
\definecolor{rojo}{HTML}{C4142C}
\definecolor{verde}{HTML}{1F9D55}
\definecolor{gris}{HTML}{5D5A73}

\newcolumntype{Y}{>{\raggedright\arraybackslash}X}
\renewcommand{\arraystretch}{1.25}
\setlist{itemsep=1pt,topsep=3pt}
\numberwithin{equation}{section}
\lstset{basicstyle=\ttfamily\small,columns=fullflexible,keepspaces=true,frame=single,rulecolor=\color{gris!40},
  literate={á}{{\'a}}1 {é}{{\'e}}1 {í}{{\'i}}1 {ó}{{\'o}}1 {ú}{{\'u}}1 {ñ}{{\~n}}1 {Ω}{{$\Omega$}}1
    {θ}{{$\theta$}}1 {λ}{{$\lambda$}}1 {Δ}{{$\Delta$}}1 {μ}{{$\mu$}}1 {·}{{$\cdot$}}1 {←}{{$\leftarrow$}}1 {§}{{\S}}1 {→}{{$\to$}}1 {≠}{{$\neq$}}1 {É}{{\'E}}1 {Á}{{\'A}}1 {∨}{{$\vee$}}1 {∧}{{$\wedge$}}1 {σ}{{$\sigma$}}1 {×}{{$\times$}}1 {≈}{{$\approx$}}1}

\newcommand{\rev}[1]{\widetilde{#1}}
\newcommand{\I}{\mathbf{I}}
\newcommand{\vect}[1]{\bm{#1}}
\newcommand{\A}{\mathcal{A}}

\newtcolorbox{meta}{colback=morado,colframe=morado,coltext=white,arc=3mm}
\newtcolorbox{lectura}{colback=morado!5,colframe=morado,fonttitle=\bfseries,title={\faBook~Lecturas}}
\newtcolorbox{ejercicio}[1]{colback=amarillo!10,colframe=amarillo!80!black,fonttitle=\bfseries,title={\faPencil*~Ejercicio #1}}
\newtcolorbox{control}{colback=rojo!5,colframe=rojo,fonttitle=\bfseries,title={\faQuestionCircle~Preguntas de control}}
\newtcolorbox{codigo}[1][De la ecuación al código]{colback=verde!6,colframe=verde,fonttitle=\bfseries,title={\faCode~#1}}
\newtcolorbox{cierre}{colback=verde!6,colframe=verde,fonttitle=\bfseries,title={\faFlagCheckered~Cierre del módulo}}
\newtcolorbox{nota}[1]{colback=amarillo!12,colframe=amarillo!80!black,fonttitle=\bfseries,title={\faLightbulb~#1}}

% \cabecera{número}{título}{subtítulo}
\newcommand{\cabecera}[3]{%
  \title{\textcolor{morado}{\textbf{Ficha #1. #2}}\\[3pt]\large Taller virtual ATA EPE · #3}
  \author{}\date{}}

\cabecera{7}{La ecuación del mecanismo}{Módulo 7: mecánica analítica II}

\begin{document}
\maketitle
\vspace{-2.5em}

\begin{meta}
\textbf{\textcolor{amarillo}{Meta del módulo.}} Para cualquier mecanismo de un grado de libertad que se arme, obtener
\emph{su} ecuación de movimiento en una sola coordenada, su inercia efectiva, su ventaja mecánica y las fuerzas que
soporta cada unión; y comprobar que coincide con la simulación completa. Cara pública: \textbf{«Ver las cuentas»} del
montaje armado.
\end{meta}

\begin{lectura}
\begin{enumerate}[leftmargin=*]
  \item \textbf{Lanczos, caps.~III--V}: principio de los trabajos virtuales, d'Alembert, ecuaciones de Lagrange y de
        Hamilton.
  \item \textbf{Goldstein, cap.~2} (§2.1--2.6) y \textbf{cap.~8} (§8.1--8.2): principio de Hamilton y hamiltoniano.
  \item \textbf{Marsden y West, «Discrete mechanics and variational integrators» (2001)}, §1--2: discretizar la
        acción en lugar de la ecuación.
  \item Arnold, \emph{Mathematical Methods of Classical Mechanics}, caps.~3--4: la mirada geométrica.
  \item Doran y Lasenby, cap.~12: lagrangianos con multivectores (opcional, para la versión en álgebra geométrica).
\end{enumerate}
\end{lectura}

% =================================================================================================
\section{Del montaje a una coordenada}

En el módulo 3, el núcleo de $\mathbf J$ dio los movimientos admisibles. En un mecanismo de un grado de libertad es
una sola dirección: para cada cuerpo $i$, su velocidad es proporcional a la de la coordenada elegida $q$ (el ángulo de
la manivela):
\begin{equation}
  \Omega_i = \xi_i(q)\,\dot q .
\end{equation}
Los bivectores $\xi_i(q)$ son los \emph{coeficientes cinemáticos}: dicen cómo se mueve cada pieza cuando la manivela
gira un poco. Se calculan normalizando el vector del núcleo para que la componente de la manivela sea 1.

% =================================================================================================
\section{El lagrangiano reducido}

Con (1.1), la energía cinética de todo el mecanismo es una forma cuadrática en $\dot q$:
\begin{equation}
  T = \tfrac12\,\mathcal M(q)\,\dot q^{\,2},\qquad \mathcal M(q) = \sum_i \xi_i(q)\vee\A_i[\xi_i(q)],\qquad
  V(q) = \sum_i m_i\,g\,y_i(q),
\end{equation}
y la ecuación de Euler–Lagrange con $\mathcal L = T - V$ da
\begin{equation}
  \mathcal M(q)\,\ddot q + \tfrac12\,\mathcal M'(q)\,\dot q^{\,2} + V'(q) = Q .
\end{equation}
$\mathcal M(q)$ es la \textbf{inercia efectiva}: lo que «siente» la mano al girar la manivela. El término con
$\mathcal M'$ hace que el mecanismo acelere o frene solo, aunque no haya fuerzas, cuando su inercia cambia con la
posición.

\begin{ejercicio}{7.1}
Deduce (2.2) desde (2.1). Para la biela-manivela de la cartilla, calcula $\mathcal M(q)$ a mano con piezas como
varillas delgadas y dibújala para $q\in[0, 2\pi]$. ¿Dónde es máxima?
\end{ejercicio}

\subsection{Fuerzas generalizadas}
Un forque $F$ (fuerza o par) aplicado al cuerpo $i$ aporta a $Q$ su potencia por unidad de $\dot q$:
\begin{equation}
  Q = \sum_k F_k\vee\xi_{i(k)}(q).
\end{equation}
El par de la mano en la manivela entra con coeficiente 1; la gravedad ya está en $V$; la fricción de las uniones,
como forques opuestos al giro relativo.

% =================================================================================================
\section{Trabajo virtual}

En equilibrio estático, $\ddot q = \dot q = 0$ y (2.2) se reduce a $V'(q) = Q$: el \textbf{principio de los trabajos
virtuales}. No hace falta conocer ninguna fuerza interna.
\begin{itemize}[leftmargin=*]
  \item La \textbf{balanza} equilibra cuando $m_1gx_1\,\delta\theta = m_2gx_2\,\delta\theta$, es decir, $m_1x_1 = m_2x_2$.
  \item La \textbf{ventaja mecánica} de cualquier transmisión ideal es la razón de desplazamientos virtuales:
        $\tau_{\text{sal}}/\tau_{\text{ent}} = \delta\theta_{\text{ent}}/\delta\theta_{\text{sal}}$, que en la
        biela-manivela cambia con el ángulo.
\end{itemize}

\begin{ejercicio}{7.2}
Calcula la ventaja mecánica de la biela-manivela en función de $q$ con los coeficientes cinemáticos. ¿Dónde es
infinita (puntos muertos) y qué significa para quien gira la manivela?
\end{ejercicio}

% =================================================================================================
\section{Las fuerzas en cada unión}

Los impulsos que el solucionador del módulo 4 calcula para cada fila son $\vect P = h\vect\lambda$: dividiendo por el
paso se obtienen los multiplicadores de Lagrange, que son las fuerzas de ligadura. En PGA, cada unión da un forque: su
parte euclídea es la fuerza que soporta el tornillo, y su parte ideal el par. El simulador las dibuja como flechas
proporcionales.

% =================================================================================================
\section{Hamilton y la conservación}

Con $p = \mathcal M(q)\dot q$, el hamiltoniano es $H = p^2/(2\mathcal M(q)) + V(q)$. Si no hay fuerzas que disipen ni
que dependan del tiempo, $H$ se conserva. Es la prueba más fina de que la simulación completa y la reducida dicen lo
mismo.

% =================================================================================================
\section{Discretización: integrador variacional}

En lugar de discretizar (2.2), se discretiza la acción $S = \int\mathcal L\,dt$ con un lagrangiano discreto
\begin{equation}
  L_d(q_k, q_{k+1}) = h\,\mathcal L\!\left(\frac{q_k + q_{k+1}}{2},\ \frac{q_{k+1}-q_k}{h}\right),
\end{equation}
y se impone que $S_d = \sum_k L_d$ sea estacionaria. Sale la ecuación de Euler–Lagrange discreta
\begin{equation}
  D_2L_d(q_{k-1}, q_k) + D_1L_d(q_k, q_{k+1}) = 0,
\end{equation}
que se resuelve para $q_{k+1}$ con el método de Newton. El método es simpléctico y conserva una versión discreta del
momento: la energía no deriva en tiempos largos.

\begin{ejercicio}{7.3}
Para $\mathcal L = \tfrac12\dot q^2 - V(q)$, muestra que (6.2) con la regla del trapecio en lugar del punto medio da el
método de Störmer–Verlet. ¿Qué cambia cuando $\mathcal M$ depende de $q$?
\end{ejercicio}

\begin{codigo}[Qué programar]
\begin{tabularx}{\linewidth}{@{}l Y@{}}
\texttt{analisis/reducido.js} & Coeficientes cinemáticos $\xi_i(q)$ desde el núcleo de $\mathbf J$; $\mathcal M(q)$,
  $\mathcal M'(q)$ por diferencias; $V(q)$; $Q$.\\
\texttt{analisis/variacional.js} & Integrador variacional (6.1)--(6.2) con Newton.\\
\texttt{analisis/trabajo\_virtual.js} & Equilibrio y ventaja mecánica.\\
\texttt{app/cuentas.js} & «Ver las cuentas»: la ecuación del mecanismo, la gráfica de $\mathcal M(q)$ y de la
  ventaja mecánica, las fuerzas en cada tornillo, la energía.\\
\end{tabularx}
\end{codigo}

\begin{control}
\begin{enumerate}[leftmargin=*]
  \item ¿Qué son los coeficientes cinemáticos y de dónde salen?
  \item ¿Por qué el mecanismo puede acelerar sin fuerza cuando $\mathcal M$ depende de $q$?
  \item ¿Por qué el trabajo virtual no necesita las fuerzas internas?
  \item ¿Qué son los puntos muertos de la biela-manivela?
  \item ¿Qué relación hay entre los impulsos del solucionador y los multiplicadores de Lagrange?
  \item ¿Qué ventaja tiene discretizar la acción en lugar de la ecuación?
\end{enumerate}
\end{control}

\begin{cierre}
Ejercicios 7.1--7.3 · pruebas P4 (balanza), P13 (trabajo virtual igual a simulación) y P14 (reducido igual a
completo) · cara pública: «Ver las cuentas» del montaje armado.
\end{cierre}

\end{document}
